\documentclass{article}
\usepackage{pathfinder-readable}

\title{Welfare Loss and Run-Level Guarantees in LLM Double Auctions}

\begin{document}
\maketitle

\begin{abstract}
This account reads a study of language-model traders in a double auction alongside a paper on calibrated
decisions in a robot's wall detector. We looked for a way to tell whether slow trading explains lost market
welfare and whether the second paper could support a guarantee for a future market run. The auction's reported
averages imply that trade count accounts for at least about 44.9\% of GPT Large's five-round welfare deficit,
but they do not identify the rest or show what caused the low count. A future-run guarantee would need
complete run records, and the ten runs reported for each model cannot give a finite 95\% bound by the proposed
method. The thread paused because it had an analytical bound and a study design, but no measured intervention
or useful market certificate.
\end{abstract}

\section{Introduction}

Struski and colleagues put language-model buyers and sellers in a continuous double auction \cite{Q}. Traders
know their own values or costs, submit prices, and can trade when an offer crosses the standing offer on the
other side. The paper reports five rounds in each of ten runs per model. Its model populations differ: GPT
Large makes relatively few trades and has low average allocative efficiency, while GPT Small comes closer to
the human benchmark. The authors link GPT Large's poor outcome partly to slow, small changes in offers that
fail to close the gap between bids and asks.

The second paper, by the Pathfinder research campaign, studies a different setting \cite{P}. A robot uses
estimated room walls to decide which scan points pass an early door-detection gate. The paper calibrates a
wall-error radius across whole, ground-truthed room episodes. Outside an uncertainty band, it guarantees that
all decisions at that gate agree with decisions made using reference walls, with a stated marginal coverage
probability. It does not guarantee the later door detections or the robot's behaviour.

The peers first asked whether that selective gate could be moved to the auction's rule for accepting a
crossing offer. The auction compares submitted prices exactly, so it has no estimated quote and reference
quote whose error needs calibration. The more useful connection was the limit of a gate guarantee: correct
trade execution says little by itself about total welfare. The peers then examined what the auction's own
numbers say about welfare loss, and how a whole-run guarantee might be posed.

\section{What was done}

The first check compared the two threshold decisions. In the robot paper, uncertainty lies in the estimated
wall. In the auction, the standing bid and ask are the prices actually posted. The peers therefore set aside a
direct price-margin transfer. They instead separated a round's lost welfare into a part forced by its number
of trades and a part due to which buyers and sellers traded. This used the auction's fixed schedule of eleven
buyer values and eleven seller costs, from 0.75 to 3.25 in steps of 0.25.

They next used the published mean trade counts to find a lower bound on the first part. A mean count does not
reveal every run's count, but counts are whole numbers. That restriction gives a stronger bound than simply
substituting the mean into a curve. They checked the calculation for GPT Large's fourth round and then for all
five rounds. Finally, they applied the robot paper's episode-level ranking method to a complete five-round
market run. They also examined a more speculative way to limit the opportunity cost of choosing one trading
pair.

\section{Results}

Let \(N\) be the number of trades in one round and \(W\) the surplus actually realised: buyer values less
seller costs, summed over the trades. Let \(f(N)\) be the greatest surplus possible with exactly \(N\) trades,
and let \(W^*\) be the greatest surplus with any number of trades. Pairing the highest-value buyers with the
lowest-cost sellers gives
\[
f(N)=2.75N-0.25N^2,\qquad W^*=7.5.
\]
The first five such trades add positive surplus; the sixth adds zero. Thus five well-chosen trades can attain
full welfare even though the paper calls six the equilibrium quantity. As the peers derived and cross-checked
in their welfare notes, every round obeys
\[
W^*-W=\bigl(W^*-f(N)\bigr)+\bigl(f(N)-W\bigr).
\]
The first term is loss forced by trade count. The second is loss from the choice and pairing of traders. Both
are non-negative. Transaction prices divide surplus between traders but do not change this sum. The identity
remains true if an offer leads an agent to trade at a loss.

The reported averages bound the count loss. For a mean count \(\mu\), let \(a\) be the whole-number part of
\(\mu\), and let \(u=\mu-a\) be its fractional part. An integer count with that mean has variance at least
\(u(1-u)\). Therefore the mean count loss is at least
\[
7.5-f(\mu)+0.25u(1-u).
\]
GPT Large's fourth round has mean count \(2.2\), giving a count loss of at least \(2.70\) surplus units. Its
reported efficiency of \(0.36\), meaning realised surplus divided by \(7.5\), implies a total mean loss of
about \(4.80\). Across all five GPT Large rounds, the minimum count losses sum to \(8.00\), against about
\(17.625\) units of reported total loss. Allowing for the efficiencies printed to two decimal places puts the
total loss at no more than \(17.8125\) units, so the count share is at least about \(44.9\%\). In round four
alone it is at least about \(55.8\%\). These bounds use the displayed counts as exact tenths, as expected for
means of ten integer counts. The peers' aggregate-welfare check gives the round-by-round arithmetic. Exact
count and selection losses still need each run's transactions.

For a possible future-run guarantee, let \(E_{ir}\) be the efficiency in round \(r\) of complete calibration
run \(i\), and let \(R_i=1-\min_{1\leq r\leq5}E_{ir}\) measure its worst round. With \(n\) exchangeable
complete runs from a frozen model, prompt, and market policy, the second paper's rank argument uses the
\(k\)th smallest score, where \(k=\lceil(n+1)(1-\alpha)\rceil\) and \(1-\alpha\) is the target coverage. It
would give a marginal lower bound for all five efficiencies in a new run. With the auction paper's \(n=10\), a
95\% target requires \(k=11\), beyond the ten available scores; the prescribed bound is therefore
uninformative. A 90\% target uses the worst observed run. The published round means do not give the
complete-run scores needed even for that calculation. This is an application of the second paper's method, not
a new conformal result.

\section{Objections and limits}

The count bound supports part of the auction paper's account of stalled trading. It does not show that small
price increments caused the low counts, and it leaves the selection loss unknown. A crossing can itself reduce
attainable welfare: the peers checked an example using buyer values \(3.25\) and \(1.00\), and seller costs
\(0.75\) and \(3.00\). Matching the high-value buyer to the cheap seller can yield \(2.50\) units. If the
low-value buyer first trades with that seller at a price of \(0.90\), a profitable crossing for both parties,
the best total after completing the remaining match is only \(0.50\). The example shows a possibility, not an
event observed in the study.

The peers also derived a conditional test for such choices. Let \(F\) be the greatest surplus attainable from
the remaining traders and \(F_e\) the greatest surplus attainable if a proposed pair \(e\) must trade. Its
opportunity cost is \(D_e=F-F_e\). If each predicted value or cost differs from the true one by at most
\(\varepsilon\), and no matching contains more than \(m\) trades, the peers' market-margin check gives \(D_e\)
at most its predicted value plus \(4m\varepsilon\). This concerns attainable surplus after a chosen pair, not
the welfare traders will actually realise. In the auction's fixed simulator the evaluator already knows all
reservation values, so it can compute this cost directly. A hidden-value version would need a fixed predictor,
varied schedules, whole-run calibration, and a score covering every state and quote history an adaptive policy
can reach. The allowance may also be too large to help. Neither paper supplies these ingredients.

The ledger notes prior work on conformal auction acceptance rules \cite{lotan2024acceptance}; the peers
therefore made no broad claim that conformal calibration in auctions is new. It also records a check of the
auction authors' code repository \cite{marketrepository}. Its README describes supplementary run files and
unseeded mechanism randomness, while the paper describes ten different random seeds. That difference matters
for a future paired-seed comparison, though it does not by itself decide whether independent runs are
exchangeable. The repository observation was recorded in the peer check and was not repeated here.

\section{Why the thread stopped}

The verifier paused the thread because its strongest quantitative finding follows from the auction paper
alone. The second paper supplies a sound way to ask about a future run, but the available runs cannot yield a
finite 95\% bound by that method. No run-level decomposition, calibrated market guarantee, or tested
intervention shows that a crossing rule improves realised welfare. The smallest next step is to use the
released transaction records to calculate both loss terms for every round of every run. A claim about
improving welfare would then need a prespecified, randomised crossing intervention, measured against the
original policy and calibrated separately for each frozen policy.

\bibliographystyle{plainurl}
\bibliography{references}

\end{document}
